% !TeX program = xelatex
% UTF-8 transcription; source wording and numerical doses are preserved.
% pic01.jpg = 20261016_220336032.jpg
% pic02.jpg = 20261016_220343175.jpg
% pic03.jpg = 20261016_220348720.jpg
% pic04.jpg = 20261016_220355193.jpg
% pic05.jpg = 20261016_220400651.jpg
% pic06.jpg = 20261016_220407612.jpg
% pic07.jpg = 20261016_220413511.jpg
% pic08.jpg = 20261016_220419875.jpg
% pic09.jpg = 20261016_220425116.jpg
% pic10.jpg = 20261016_220431562.jpg
% pic11.jpg = 20261016_220437143.jpg
\newif\ifPNSStandalone
\ifdefined\settextfont\PNSStandalonefalse\else\PNSStandalonetrue\fi
\ifPNSStandalone
  \documentclass[a4paper,12pt]{article}
  \usepackage[margin=2.2cm]{geometry}
  \usepackage{amsmath,amssymb}
  \usepackage{enumitem}
  \usepackage{xepersian}
  \settextfont{Amiri}
  \setlatintextfont{FreeSerif}
\fi
\ifPNSStandalone\begin{document}\fi
\begingroup
\emergencystretch=3em
\newcommand{\PNSen}[1]{\lr{#1}}
\newcommand{\PNSunclear}[1]{[خوانش نامطمئن: #1]}
\newcommand{\PNSpage}[2]{\clearpage\section*{صفحهٔ #1}\noindent\PNSen{\texttt{images/pic#2.jpg}}\par\smallskip
{\footnotesize\PNSen{Note book \quad Subject: 93 \quad Veor: \quad Month: \quad Date:}}\par\medskip}
\newcommand{\PNSfooter}[1]{\par\medskip{\footnotesize\PNSen{BAHARE DANESH ASL}\hfill #1}\par}

\PNSpage{۱}{01}
\noindent\PNSen{Subject: Harrison -- 93}
\section*{\PNSen{Paraneoplastic Syndrome: Endocrinologic/Hematologic}}
\begin{enumerate}[label=\arabic*.,start=1]
\item تومورهایی با منشأ نورواندوکرین علل شایع سندرم‌های پارانئوپلاستیک هستند:
\begin{itemize}
\item \PNSen{Small Cell Lung Carcinoma (SCLC)}
\item \PNSen{Carcinoids}
\end{itemize}
\item بسیاری از موارد سندرم پارانئوپلاستیک تشخیص داده می‌شوند.
\begin{itemize}
\item هرگونه تظاهر بالینی غیرمعمول در بیمار سرطانی باید ما را مشکوک به سندرم‌های پارانئوپلاستیک نماید.
\end{itemize}
\item منشأ تولید هورمون‌ها: \PNSen{Eutopic} / \PNSen{Ectopic}.
\item بسیاری از هورمون‌ها به‌طور معمول به میزان بسیار اندک توسط بقیهٔ بافت‌ها غیر از بافت مورد انتظار نیز تولید می‌شوند.
\item در تولید اکتوپیک هورمون:
\begin{itemize}
\item میزان آن هورمون بالا می‌رود.
\item فیدبک تنظیمی آن هورمون به هم می‌خورد.
\item با پروسسینگ دقیق، معمولاً پرکرسورهای آن هورمون.
\end{itemize}
\item مکانیسم‌های مولکولی تولید اکتوپیک هورمون:
\begin{latin}
\begin{itemize}
\item[a)] Genetic rearrangement (rare)
\begin{itemize}
\item Translocation of PTH gene
\item In leukemia and lymphoma:
\[\text{Somatic genetic rearrangement}\longrightarrow\text{Alter cellular differentiation}\]
\end{itemize}
\end{itemize}
\end{latin}
(تولید \PNSen{PTH} در بافت‌های غیر پاراتیروئید.)
\end{enumerate}
\PNSfooter{۱}

\PNSpage{۲}{02}
\begin{latin}
\noindent b) Excessive production of numerous peptides.
\end{latin}
\begin{enumerate}[label=\arabic*.,start=7]
\item احتمالاً تمایز سلولی علت اکثر موارد تولید نابجای هورمون است.
\item \PNSen{Many cancers are poorly differentiated.}
\begin{latin}
\noindent Certain tumor products such as:
\begin{itemize}
\item PTH related pro (PTHrP)
\item hCG
\item $\alpha$ fetoprotein
\end{itemize}
Are characteristic of gene expression at earlier developmental stages.
\end{latin}
\item در \PNSen{SCLC} سلول‌ها نسبتاً تمایزنیافته هستند.
\item شایع‌ترین سندرم‌های پارانئوپلاستیک اندوکرین:
\begin{latin}
\begin{enumerate}[label=\alph*)]
\item Hypercalcemia from overproduction of PTHrP
\item Hyponatremia from excess vasopressin
\item Cushing's syndrome from ectopic ACTH
\end{enumerate}
\end{latin}
\end{enumerate}
\section*{هایپرکلسمی ناشی از \PNSen{PTHrP} اکتوپیک}
\begin{latin}\noindent Humoral Hypercalcemia of Malignancy\end{latin}
\begin{enumerate}[label=\arabic*.,start=1]
\item \PNSen{HHM} در ۲۰٪ بیماران سرطانی رخ می‌دهد.
\item شایع‌ترین بدخیمی‌های عامل \PNSen{HHM}: ریه / سروگردن / پوست / مری / پستان / ادراری تناسلی / \PNSen{MM} / لنفوم.
\end{enumerate}
\PNSfooter{۲}

\PNSpage{۳}{03}
\begin{enumerate}[label=\arabic*.,start=3]
\item \PNSen{HHM} علل مختلفی دارد ولی شایع‌ترین آن تولید بیش از حد \PNSen{PTHrP} است.
\item \PNSen{PTHrP may affect the initiation / progression of tumors by acting through pro-survival and chemokine pathways.}
\item \PNSen{PTHrP} شبیه به \PNSen{PTH} است و به رسپتور آن متصل می‌شود.
\item \PNSen{PTHrP} نقش فیزیولوژیک مهمی در موارد زیر دارد:
\begin{latin}
\begin{itemize}
\item Skeletal development
\item Cellular proliferation regulation
\item Differentiation in other tissues
\end{itemize}
\end{latin}
شامل پوست / مغز استخوان / پستان / فولیکول مو.
\item \PNSen{TGF-$\beta$} که توسط بسیاری از تومورها تولید می‌شود، \PNSen{PTHrP} را فعال می‌کند.
\item موتاسیون ژن \PNSen{Ras} بیان \PNSen{PTHrP} را افزایش می‌دهد.
\item از دست دادن ژن سرکوبگر تومور \PNSen{p53} $\leftarrow$ افزایش تولید \PNSen{PTHrP}.
\item تولید بیش از حد \PNSen{1,25 vit D} $\leftarrow$ علت شایع دیگر \PNSen{HHM}.
\begin{itemize}
\item \PNSen{Lymphomas and granulomatous disorders if associated with Ca$\uparrow$}
\item آنزیم نابجا که \PNSen{25 vit D} را تبدیل به \PNSen{1,25 vit D} می‌کند.
\item افزایش جذب روده‌ای کلسیم.
\end{itemize}
\end{enumerate}
\PNSfooter{۳}

\PNSpage{۴}{04}
\begin{enumerate}[label=\arabic*.,start=11]
\item علل دیگر \PNSen{HHM}: تولید سایتوکاین‌های استئولیتیک و فاکتورهای التهابی توسط برخی تومورها.
\item تظاهرات بالینی \PNSen{HHM}:
\begin{itemize}
\item بیماری با یک بدخیمی شناخته‌شده که هایپرکلسمیک می‌شود.
\item با احتمال کمتر هایپرکلسمی ممکن است تظاهر اولیهٔ بدخیمی باشد.
\end{itemize}
\item وقتی سطح کلسیم از \PNSen{14 mg/dl} بالاتر برود:
\begin{latin}
\begin{itemize}
\item Fatigue
\item Mental status changes
\item Polyuria
\item Dehydration
\item Symptoms of nephrolithiasis
\end{itemize}
\end{latin}
\item عوارض قلبی هایپرکلسمی:
\begin{latin}
\begin{itemize}
\item Shorten ST segment
\item Shorten QT interval
\item Bundle branch blocks
\end{itemize}
\end{latin}
و برادی‌آریتمی‌ها.
\item تشخیص \PNSen{HHM}:
\begin{itemize}
\item بدخیمی شناخته‌شده
\item هایپرکلسمی \PNSen{Recent onset}
\item سطوح بسیار بالای کلسیم خون
\end{itemize}
\end{enumerate}
\PNSfooter{۴}

\PNSpage{۵}{05}
\begin{enumerate}[label=\arabic*.,start=16]
\item مانند هایپرپاراتیروئیدی، \PNSen{HHM} نیز منجر به هایپرکلسی‌اوری و هایپوفسفاتمی می‌شود.
\item در \PNSen{HHM}: آلکالوز متابولیک / در هایپرپارا: اسیدوز هایپرکلرمیک.
\item اندازه‌گیری سطح خونی \PNSen{PTH} برای رد هایپرپارا اولیه.
\item تأیید تشخیص \PNSen{HHM}: سطح بالای \PNSen{PTHrP}؛ در ۸۰٪ بیماران هایپرکلسمی با سرطان بالا است.
\item در بیماران لنفوم مبتلا به هایپرکلسمی ممکن است سطح \PNSen{1,25 vit D} بالا باشد.
\end{enumerate}
\section*{درمان \PNSen{HHM}}
\begin{latin}
\begin{enumerate}[label=\arabic*.,start=1]
\item Removal of excess Ca in:
\begin{itemize}
\item Diet
\item Medications
\item IV solutions
\end{itemize}
\item Saline rehydration (200--500 mL/h)
\begin{itemize}
\item To dilute serum calcium
\item To promote calciuresis
\end{itemize}
\item Exercise caution in patients with cardiac / hepatic / renal insufficiency.
\item Forced diuresis with:
\begin{itemize}
\item Furosemide (20--80 mg IV, escalating doses)
\item Other loop diuretics
\end{itemize}
\begin{itemize}
\item Only in life-threatening hypercalcemia
\item Only after complete rehydration
\end{itemize}
\end{enumerate}
\end{latin}
\PNSfooter{۵}

\PNSpage{۶}{06}
\begin{latin}
\begin{enumerate}[label=\arabic*.,start=5]
\item Oral phosphorus (250 mg Neutra-Phos 3--4 times daily)
\begin{itemize}
\item Until serum P $>3$ mg/dl
\end{itemize}
\item Bisphosphonates (Pamidronate, Zoledronate) $\longrightarrow$ For chronic treatment
\end{enumerate}
\end{latin}
کلسیم سرم را ظرف ۱ تا ۴ روز کاهش می‌دهند و آزاد شدن کلسیم را برای چند هفته سرکوب می‌کنند.
\begin{latin}
\begin{enumerate}[label=\arabic*.,start=7]
\item Denosumab (120 mg SC weekly for 4 weeks then monthly)
\begin{itemize}
\item In patients who do not respond adequately to bisphosphonates.
\item Decoy receptor for RANK ligand to mitigate stimulation of osteoclasts.
\end{itemize}
\item Cinacalcet (30 mg PO bid to 90 mg PO qid)
\begin{itemize}
\item Stimulates calcium-sensing receptors to suppress PTH secretion.
\end{itemize}
\end{enumerate}
\end{latin}
برای کارسینوم پاراتیروئید یا موارد نادر ترشح نابجای \PNSen{PTH} در اثر تومور.
\begin{latin}
\begin{enumerate}[label=\arabic*.,start=9]
\item Glucocorticoid (Prednisone 40--100 mg PO in four divided doses)
\begin{itemize}
\item In hypercalcemia associated with lymphomas / MM / leukemia
\end{itemize}
\item Dialysis $\longrightarrow$ In severe hypercalcemia
\end{enumerate}
\end{latin}
* قبلاً از کلسی‌تونین و میترامایسین استفاده می‌شد که خیلی سودی ندارند.
\PNSfooter{۶}

\PNSpage{۷}{07}
\section*{\PNSen{Tumor-Associated Syndrome of Inappropriate Antidiuretic Hormone (SIADH)}}
\begin{enumerate}[label=\arabic*.,start=1]
\item \PNSen{SIADH} در نیمی از بیماران \PNSen{SCLC} رخ می‌دهد.
\item \PNSen{SIADH} می‌تواند در اثر ترومای \PNSen{CNS}، عفونت‌ها و داروها نیز رخ دهد.
\item در واکنش‌های \PNSen{SIADH} به‌طور جبرانی تشنگی کاهش می‌یابد و جلوی ایجاد هیپوناترمی به‌طور موقت گرفته می‌شود.
\item در \PNSen{SIADH} نابجا، مایع خوراکی یا وریدی به‌سرعت منجر به هیپوناترمی می‌شود چون دیورز کلیوی کاهش یافته است.
\item تومورهایی با منشأ نورواندوکرین مثل \PNSen{SCLC} / کارسینوئیدها مهم‌ترین علل تولید نابجای وازوپرسین.
\item بقیهٔ بدخیمی‌ها به‌ندرت نیز می‌توانند منجر به تولید نابجای وازوپرسین شوند: بدخیمی‌های تخمدان، \PNSen{GI}، ادراری تناسلی و هموپویتیک.
\item مکانیسم فعال شدن ژن وازوپرسین در این تومورها ناشناخته است.
\end{enumerate}
\section*{تظاهرات بالینی}
\begin{enumerate}[label=\arabic*.,start=8]
\item اغلب بیماران با تولید نابجای وازوپرسین بدون علامت هستند.
\begin{itemize}
\item در یک آزمایش روتین متوجه هیپوناترمی می‌شوند.
\end{itemize}
\end{enumerate}
\PNSfooter{۷}

\PNSpage{۸}{08}
\begin{enumerate}[label=\arabic*.,start=9]
\item علائم شامل:
\begin{latin}
Weakness, Lethargy, Nausea, Confusion, Depressed mental status, Seizure.
\end{latin}
\item شدت علائم وابسته به:
\begin{itemize}
\item شدت هیپوناترمی
\item سرعت ایجاد هیپوناترمی
\item هیپوناترمی معمولاً آهسته به وجود می‌آید مگر آنکه بیمار مایع وریدی دریافت کند.
\end{itemize}
\item تشخیص:
\begin{itemize}
\item مانند بقیهٔ علل \PNSen{SIADH} است.
\item هیپوناترمی
\item کاهش اسمولالیتهٔ سرم $\leftarrow$ اسمولالیتهٔ ادرار افزایش یافته است.
\item سطح سدیم در ادرار افزایش می‌یابد.
\item رد سایر علل هیپوناترمی $\leftarrow$ نارسایی کلیه، آدرنال و هیپوتیروئیدی.
\end{itemize}
\end{enumerate}
* اندازه‌گیری سطح وازوپرسین عموماً برای تشخیص ضروری نیست.
\section*{درمان}
\begin{enumerate}[label=\arabic*.,start=12]
\item بیماران با وازوپرسین نابجا ظرف چند هفته تا چند ماه مبتلا به هیپوناترمی می‌شوند.
\item سطح سدیم باید تدریجی اصلاح شود مگر:
\begin{itemize}
\item تغییر وضعیت هوشیاری رخ داده باشد.
\item ریسک تشنج وجود داشته باشد.
\end{itemize}
\item اصلاح سریع سدیم $\leftarrow$ دمیلیناسیون مغز $\leftarrow$ \PNSen{Central pontine myelinolysis}.
\end{enumerate}
\PNSfooter{۸}

\PNSpage{۹}{09}
\begin{enumerate}[label=\arabic*.,start=15]
\item درمان بدخیمی زمینه‌ای ممکن است کمک‌کننده باشد، ولی زمان‌بر است.
\item \PNSen{Fluid restriction to less than urine output, plus insensible losses.}
\begin{itemize}
\item عمدتاً برای اصلاح هیپوناترمی کافی است.
\end{itemize}
\item سرم سالین و قرص‌های نمک کمک‌کننده نیستند مگر:
\begin{latin}\noindent Volume depletion is also present.\end{latin}
\item \PNSen{Demeclocycline (150--300 mg orally 3--4 times daily)}
\begin{latin}
\begin{itemize}
\item Inhibit vasopressin action on the renal distal tubule
\item Onset of action is slow (1--2 weeks)
\end{itemize}
\end{latin}
\item \PNSen{Vaptan class of drugs}
\begin{latin}
\noindent Inhibiting vasopressin receptors in the renal collecting ducts.
\begin{itemize}
\item Conivaptan $\longrightarrow$ $V_2$-receptor antagonist
\begin{itemize}
\item Effective when used in combination with fluid restriction.
\item PO and IV
\end{itemize}
\item Tolvaptan (15 mg PO daily) $\longrightarrow$ Increased to 30--60 mg daily
\end{itemize}
\end{latin}
\item \PNSen{Hypertonic (3\%) or N/S + furosemide}
\begin{latin}
\begin{itemize}
\item Severe hyponatremia (Na $<115$ meq/L) or mental status changes
\end{itemize}
\end{latin}
\item \PNSen{Rate of Na correction (0.5--1 meq/L per hour)} *
\end{enumerate}
\PNSfooter{۹}

\PNSpage{۱۰}{10}
\section*{سندرم کوشینگ ناشی از \PNSen{ACTH} نابجا}
\begin{enumerate}[label=\arabic*.,start=1]
\item ۵ تا ۱۰٪ موارد سندرم کوشینگ را شامل می‌شود.
\item شایع در تومورهای نورواندوکرین:
\begin{itemize}
\item \PNSen{SCLC}: شایع‌ترین علت
\item کارسینوئید برونکیال و \PNSen{thymic}
\item \PNSen{Islet cell tumors}
\item فئوکروموسایتوما
\end{itemize}
\item چون تومورها فاقد بسیاری از آنزیم‌های مورد نیاز در مسیر تولید \PNSen{ACTH} هستند، بخش کمی از موارد تولیدشده به‌صورت \PNSen{ACTH} فعال است و عمدهٔ آن پروتئین‌های بزرگ شکسته‌نشدهٔ غیرفعال هستند.
\item برخی تومورها می‌توانند \PNSen{CRH} اکتوپیک تولید کنند:
\begin{latin}
\begin{itemize}
\item SCLC
\item Pancreatic islet cell tumors
\item MTC
\item Carcinoids
\item Prostate cancer
\end{itemize}
\end{latin}
\begin{itemize}
\item زمانی که سطح \PNSen{CRH} به اندازهٔ کافی بالا برود، منجر به سندرم کوشینگ می‌شود.
\item گاهی اوقات تومورها هم \PNSen{CRH} و هم \PNSen{ACTH} تولید می‌کنند.
\end{itemize}
\item \PNSen{ACTH-independent Cushing's syndrome}: در سلول‌های آدرنال.
\begin{latin}
\noindent Ectopic expression of the gastric inhibitory peptide (GIP) receptor.
\end{latin}
خوردن غذا $\leftarrow$ تحریک ترشح \PNSen{GIP} $\leftarrow$ تحریک غدهٔ آدرنال $\leftarrow$ تولید گلوکوکورتیکوئید.
\end{enumerate}
\PNSfooter{۱۰}
\noindent{\footnotesize حاشیهٔ پایین تصویر: ۱۱؛ \PNSen{Severe hypokalemia}.}

\PNSpage{۱۱}{11}
\begin{enumerate}[label=\arabic*.,start=6]
\item تظاهرات بالینی:
\begin{itemize}
\item صرفاً کمی از بیماران با تولید \PNSen{ACTH} نابجا علائم بالینی هایپرکورتیزولیسم را پیدا می‌کنند.
\end{itemize}
\begin{latin}
\noindent Patients with ectopic ACTH:
\begin{itemize}
\item Less marked weight gain
\item Less centripetal fat redistribution
\end{itemize}
\end{latin}
به علت برخورد کوتاه‌مدت به کورتیزول و کاشکسی ناشی از بدخیمی.
\begin{latin}
\noindent Metabolic manifestations of ectopic ACTH syndrome:
\begin{itemize}
\item Fluid retention
\item HTN
\item Hypokalemia
\item Metabolic alkalosis
\item Glucose intolerance
\item Steroid psychosis
\end{itemize}
\end{latin}
\begin{itemize}
\item سطوح بالای \PNSen{ACTH} منجر به افزایش پیگمانتاسیون می‌شود.
\item سطح بسیار بالای گلوکوکورتیکوئید:
\begin{latin}
\begin{itemize}
\item Marked skin fragility
\item Easy bruising
\end{itemize}
\end{latin}
\item سطح بالای کورتیزول $\leftarrow$ آنزیم \PNSen{$11\beta$-hydroxysteroid dehydrogenase II} نتواند جلوی اتصال کورتیزول به گیرندهٔ \PNSen{Renal mineralocorticoid receptors} را بگیرد $\leftarrow$ \PNSen{Severe hypokalemia}.
\end{itemize}
\end{enumerate}
\PNSfooter{۱۱}
\endgroup
\ifPNSStandalone
  \def\PNSfinish{\end{document}}
\else
  \let\PNSfinish\relax
\fi
\PNSfinish
